\section{Remove K Digits}
\label{sec:sq-remove-k-digits}

\problemstatement{Given a non-negative integer represented as a string
$\textit{num}$, and an integer $k$, remove exactly $k$ digits from
$\textit{num}$ so that the remaining digits, kept in their original
relative order, form the \textbf{smallest possible} number. The result
must not have leading zeros (unless it is exactly \texttt{"0"}).}

\begin{patternbox}
\emph{``Remove $k$ elements (keeping relative order) to make the result as
small/large as possible''} is a greedy-with-a-stack pattern: a number is
smaller when an earlier (more significant) digit is smaller, so whenever
we see a digit smaller than what we have just placed, removing that
just-placed (now regretted) larger digit is always an improvement -- a
monotonic increasing stack of kept digits, with removals funded by the
budget $k$, captures exactly this.
\end{patternbox}

\subsection*{Understanding the problem carefully}

Among any two digit strings of the same length, the one with the smaller
digit at the first position where they differ is the smaller number --
so, greedily, whenever we have the opportunity to remove a digit and
thereby make some earlier-kept digit smaller (by exposing a smaller digit
to take its place at that position), we should take it, as long as we
still have removals left in our budget $k$.

\begin{ideabox}[Monotonic Increasing Stack of Kept Digits, Funded by Budget $k$]
Scan $\textit{num}$ left to right, maintaining a stack of digits kept so
far. For each new digit $d$: while $k>0$ and the stack is non-empty and
its top digit exceeds $d$, pop the top (spend one unit of $k$) -- this
digit was a mistake to keep, now that a smaller one is available to take
its place. Then push $d$. After the scan, if $k$ still has unused budget,
remove that many digits from the \emph{end} of the stack (the largest
remaining suffix, since no smaller digit ever showed up to justify
removing them earlier). Finally, strip any leading zeros, and return
\texttt{"0"} if the result is empty.
\end{ideabox}

\subsection*{Why popping a larger top whenever a smaller digit arrives is optimal}

Removing a digit only ever helps if it lets an earlier, larger digit be
replaced (in significance/position) by a smaller later digit -- and the
earliest opportunity to do this should always be taken, since a smaller
digit at a \emph{more significant} (more leftward) position always
beats any improvement achievable further right. This is exactly the same
exchange-argument logic as every monotonic-stack pop rule in this
chapter: keeping a larger digit when a smaller one is available and
removals remain in budget can never be better than removing it, since the
resulting number would be lexicographically (and numerically, for
same-length strings) larger or equal in every case.

\subsection*{Algorithm}

\begin{enumerate}
  \item Initialize an empty stack of digits.
  \item For each digit $d$ in $\textit{num}$:
  \begin{enumerate}
    \item While $k>0$ and the stack is non-empty and its top $>d$: pop;
          $k \gets k-1$.
    \item Push $d$.
  \end{enumerate}
  \item While $k>0$: pop from the end of the stack; $k \gets k-1$.
  \item Strip leading zeros from the stack's contents (read bottom to
        top). If the result is empty, return \texttt{"0"}.
\end{enumerate}

\subsection*{Dry run}

\dryrun{Take $\textit{num} = \texttt{"1432219"}$, $k=3$.}

\begin{itemize}
  \item \texttt{1}: stack empty, push. Stack $=[1]$.
  \item \texttt{4}: top $1<4$, no pop. Push. Stack $=[1,4]$.
  \item \texttt{3}: top $4>3$ and $k=3>0$: pop $4$, $k=2$. New top
        $1<3$, stop. Push $3$. Stack $=[1,3]$.
  \item \texttt{2}: top $3>2$ and $k=2>0$: pop $3$, $k=1$. New top
        $1<2$, stop. Push $2$. Stack $=[1,2]$.
  \item \texttt{2}: top $2 \not> 2$, no pop. Push. Stack $=[1,2,2]$.
  \item \texttt{1}: top $2>1$ and $k=1>0$: pop $2$, $k=0$. New top
        $2>1$ but $k=0$ now, stop. Push $1$. Stack $=[1,2,1]$.
  \item \texttt{9}: $k=0$, no pop. Push. Stack $=[1,2,1,9]$.
\end{itemize}

$k=0$, no trailing removal needed. Result (no leading zeros to strip):
\texttt{"1219"}.

\subsection*{C++ implementation}

\begin{lstlisting}
#include <bits/stdc++.h>
using namespace std;

string removeKdigits(string num, int k) {
    string st; // used as a stack of characters

    for (char d : num) {
        while (k > 0 && !st.empty() && st.back() > d) {
            st.pop_back();
            k--;
        }
        st.push_back(d);
    }

    while (k > 0 && !st.empty()) {
        st.pop_back();
        k--;
    }

    int i = 0;
    while (i < (int)st.size() - 1 && st[i] == '0') i++;
    st = st.substr(i);

    return st.empty() ? "0" : st;
}

int main() {
    cout << removeKdigits("1432219", 3) << endl;
    return 0;
}
\end{lstlisting}

\begin{complexitybox}
\textbf{Time Complexity.} Each digit is pushed once and popped at most
once, giving
\[
  O(n).
\]
\textbf{Space Complexity.} $O(n)$ for the stack.
\end{complexitybox}

\begin{takeawaybox}
``Remove exactly $k$ elements to minimize/maximize the resulting
sequence, preserving relative order'' is solved by a monotonic stack
where the popping budget is explicitly $k$ -- once $k$ reaches zero, the
stack stops discarding even objectively ``bad'' elements, and any
leftover budget at the end is spent removing from the tail, since nothing
better ever arrived to justify earlier removals there.
\end{takeawaybox}
